% =====================================================================
%  Section: Architecture
%
%  Destination is the ISFPGA_Submission repository, which is a separate
%  repo synced to Overleaf and gitignored here. This copy is staged
%  beside the figure it includes so the two stay in step; copy both into
%  the paper project, or \input this file from main.tex once the figure
%  is on its TEXINPUTS path.
%
%  PROSE IS NOT YET IN THE PROJECT VOICE. CLAUDE.md requires reading
%  ISFPGA_Submission/persona.md before this reaches main.tex, and that
%  file is not reachable from this repository. Treat the wording as a
%  content draft.
%
%  Preamble needs:
%      \usepackage{tikz}
%      \usetikzlibrary{arrows.meta}
%      \usepackage{amsmath}
%
%  Symbols used, all introduced in the text on first use:
%      H     KV heads per layer                        (8)
%      G     query heads sharing one KV head           (4)
%      d     head dimension                            (64)
%      b_k   bits per key codeword                     (4)
%      b_v   bits per value codeword                   (2)
%      w_n   bits per stored norm                      (16)
%      R     bits in one cache row                     (416, i.e. 52 B)
%      P     slices the row is cut into, one per port  (4)
%      T     tokens the scan covers
% =====================================================================

\begin{figure*}[t]
  \centering
  % =====================================================================
%  Figure: the attention block on the ZCU104
%
%  Include from the paper:
%
%      \begin{figure*}[t]
%        \centering
%        % =====================================================================
%  Figure: the attention block on the ZCU104
%
%  Include from the paper:
%
%      \begin{figure*}[t]
%        \centering
%        % =====================================================================
%  Figure: the attention block on the ZCU104
%
%  Include from the paper:
%
%      \begin{figure*}[t]
%        \centering
%        \input{fig_hardware}
%        \caption{...}\label{fig:hw}
%      \end{figure*}
%
%  Preamble needs:
%      \usepackage{tikz}
%      \usetikzlibrary{arrows.meta}
%
%  NOT YET COMPILED -- written without a LaTeX toolchain available.
%  Check it renders before trusting it.
%
%  ---------------------------------------------------------------------
%  Geometry is in the same 1280 x 455 coordinate space as
%  fig_hardware.svg, so the two can be compared line for line.  The unit
%  vectors below set the printed size; text is set in real LaTeX sizes
%  and deliberately does NOT scale with the drawing.
%
%  To resize the drawing   : change \figunit
%  To resize the type      : change \figTitle / \figMod / \figNote
%
%  \figunit = 0.13906mm gives a width of 178mm (7in), a two-column
%  figure*.  For a single-column figure use about 0.0664mm, but at that
%  size the type will have to come down with it.
%
%  TIGHTEST LABEL: \texttt{softmax\_online} in a 138-unit box.  At
%  \figunit = 0.13906mm that box is 54.4pt wide, and the label needs
%  roughly 59pt at \scriptsize.  If it overflows, either set \figMod to
%  \tiny or widen that column.
% =====================================================================

\providecommand{\figunit}{0.13906mm}
\providecommand{\figTitle}{\small}      % Host, DDR
\providecommand{\figMod}{\scriptsize}   % module names, section headings
\providecommand{\figNote}{\tiny}        % annotations and body text

\definecolor{figink}{HTML}{1A1A1A}
\definecolor{fighost}{HTML}{F2F2F2}
\definecolor{figddr}{HTML}{D8D8D8}
\definecolor{figlane}{HTML}{FAFAFA}
\definecolor{figcell}{HTML}{F0F0F0}
\definecolor{figrulec}{HTML}{9A9A9A}
\definecolor{figdim}{HTML}{555555}
\definecolor{figbody}{HTML}{333333}
\definecolor{figleg}{HTML}{666666}

\begin{tikzpicture}[
    x = {(\figunit,0mm)},
    y = {(0mm,\figunit)},
    line width = 0.5pt,
    draw       = figink,
    every node/.style = {inner sep=0pt, outer sep=0pt, text=figink},
    box/.style    = {draw=figink, fill=white},
    arw/.style    = {draw=figink, -{Stealth[length=3.2pt,width=2.6pt]}},
    rul/.style    = {draw=figrulec, line width=0.3pt},
    mod/.style    = {anchor=base, font=\figMod\ttfamily},
    ttl/.style    = {anchor=base, font=\figTitle},
    hed/.style    = {anchor=base west, font=\figMod},
    nte/.style    = {anchor=base, font=\figNote, text=figbody},
    dim/.style    = {anchor=base, font=\figNote, text=figdim},
    nteW/.style   = {anchor=base west, font=\figNote, text=figbody},
    dimW/.style   = {anchor=base west, font=\figNote, text=figdim},
]

% ---------------------------------------------------------------- regions
\draw[box, fill=fighost] (40,35)   rectangle (205,345);
\draw[box]               (235,35)  rectangle (1045,345);
\draw[box, fill=figddr]  (1075,35) rectangle (1250,345);

% ------------------------------------------------- write path, left to right
\draw[box] (250,222) rectangle (388,285);
\draw[box] (410,222) rectangle (548,285);
\draw[box] (570,222) rectangle (708,285);
\draw[box] (730,222) rectangle (868,285);
\draw[box] (890,222) rectangle (1028,285);

% --------------------- lane array; kv_store_ddr sits outside it, being shared
\draw[box, fill=figlane] (240,75) rectangle (872,163);
\draw[box] (250,88) rectangle (388,151);
\draw[box] (410,88) rectangle (548,151);
\draw[box] (570,88) rectangle (708,151);
\draw[box] (730,88) rectangle (868,151);

\draw[box] (890,77) rectangle (1028,161);
\draw[box, fill=figcell] (897,89) rectangle (927,117);
\draw[box, fill=figcell] (933,89) rectangle (963,117);
\draw[box, fill=figcell] (987,89) rectangle (1021,117);

% ---------------------------------------------------------------- arrows
\draw[arw] (208,254)  -- (246,254);
\draw[arw] (390,254)  -- (408,254);
\draw[arw] (550,254)  -- (568,254);
\draw[arw] (710,254)  -- (728,254);
\draw[arw] (870,254)  -- (888,254);
\draw[arw] (1032,254) -- (1071,254);

\draw[arw] (1071,119) -- (1032,119);
\draw[arw] (888,119)  -- (869,119);
\draw[arw] (728,119)  -- (710,119);
\draw[arw] (568,119)  -- (550,119);
\draw[arw] (408,119)  -- (390,119);
\draw[arw] (246,119)  -- (208,119);

\draw[arw] (479,222) -- (479,191) -- (799,191) -- (799,152);

% ------------------------------------------------- rules in host and DDR
\draw[rul] (56,301)   -- (189,301);
\draw[rul] (56,197)   -- (189,197);
\draw[rul] (56,123)   -- (189,123);
\draw[rul] (1090,285) -- (1235,285);
\draw[rul] (1090,235) -- (1235,235);
\draw[rul] (1090,143) -- (1235,143);

% ---------------------------------------------------------------- host
\node[ttl]  at (122,317) {Host};
\node[hed]  at (56,279)  {sent to the block};
\node[dimW] at (64,260)  {$q$, $k$ and $v$ for every};
\node[dimW] at (64,245)  {head, already projected};
\node[dimW] at (64,230)  {through $W_q$, $W_k$, $W_v$};
\node[dimW] at (64,215)  {and carrying RoPE};
\node[hed]  at (56,175)  {back from the block};
\node[dimW] at (64,156)  {the attention output};
\node[dimW] at (64,141)  {for all of the heads};
\node[hed]  at (56,101)  {finished on the host};
\node[dimW] at (64,82)   {$W_o'$ \ ($R^{-1}$ folded in)};
\node[dimW] at (64,67)   {MLP \ $\rightarrow$ \ logits};
\node[dimW] at (64,52)   {then the next layer};

% ---------------------------------------------------------------- DDR
\node[ttl] at (1162,319) {DDR};
\node[anchor=base, font=\figMod] at (1162,299) {KV cache};
\node[nte] at (1162,265) {one 52-byte row per};
\node[nte] at (1162,251) {token, per KV head};
\node[nte] at (1162,215) {the row is cut into $P$};
\node[nte] at (1162,201) {slices, and each one is};
\node[nte] at (1162,187) {kept as its own array,};
\node[nte] at (1162,173) {contiguous and running};
\node[nte] at (1162,159) {across all the tokens};
\node[nte] at (1162,123) {each AXI port reads};
\node[nte] at (1162,109) {one slice and nothing};
\node[nte] at (1162,95)  {else, which turns the};
\node[nte] at (1162,81)  {whole scan into long};
\node[nte] at (1162,67)  {sequential bursts};

% ------------------------------------------------------- programmable logic
\node[hed]  at (250,322) {Programmable logic};
\node[dimW] at (250,305) {one KV head at a time --- the $H$ heads of a token run one after the other};
\node[nteW] at (240,169) {$\times\,G$ \ --- one lane per query head of this KV head};

% ---------------------------------------------------------------- write path
\node[mod] at (319,249) {attn\_ingress};
\node[mod] at (479,258) {rot\_fwht};
\node[dim] at (479,239) {$k$, $v$, $q$};
\node[mod] at (639,249) {rot\_norm};
\node[mod] at (799,249) {rot\_encode};
\node[mod] at (959,262) {kv\_write};
\node[dim] at (959,245) {one row per token,};
\node[dim] at (959,233) {per KV head};

% ---------------------------------------------------------------- lanes
\node[mod] at (319,115) {finalize};
\node[mod] at (479,115) {accum};
\node[mod] at (639,115) {softmax\_online};
\node[mod] at (799,115) {score\_lane};

% ------------------------------------------------------- shared read engine
\node[mod] at (959,138) {kv\_store\_ddr};
\node[nte] at (912,98)  {0};
\node[nte] at (948,98)  {1};
\node[dim] at (975,97)  {$\dots$};
\node[nte] at (1004,98) {$P\!-\!1$};
\node[nte] at (959,64)  {the $P$ ports fetch one};
\node[nte] at (959,52)  {whole row per cycle, and};
\node[nte] at (959,40)  {all $G$ lanes score on it};

% ---------------------------------------------------------------- q branch
\node[nte] at (639,198) {$q$ is held inside each lane and never written to the cache};

% ---------------------------------------------------------------- legend
\node[anchor=base west, font=\figNote, text=figleg] at (40,15)
  {The ZCU104 provides four AXI-HP ports, each carrying 16 bytes a cycle. A plane never straddles a field boundary, so the key codes, the value codes and the norms must together fall into at most four planes};

\end{tikzpicture}

%        \caption{...}\label{fig:hw}
%      \end{figure*}
%
%  Preamble needs:
%      \usepackage{tikz}
%      \usetikzlibrary{arrows.meta}
%
%  NOT YET COMPILED -- written without a LaTeX toolchain available.
%  Check it renders before trusting it.
%
%  ---------------------------------------------------------------------
%  Geometry is in the same 1280 x 455 coordinate space as
%  fig_hardware.svg, so the two can be compared line for line.  The unit
%  vectors below set the printed size; text is set in real LaTeX sizes
%  and deliberately does NOT scale with the drawing.
%
%  To resize the drawing   : change \figunit
%  To resize the type      : change \figTitle / \figMod / \figNote
%
%  \figunit = 0.13906mm gives a width of 178mm (7in), a two-column
%  figure*.  For a single-column figure use about 0.0664mm, but at that
%  size the type will have to come down with it.
%
%  TIGHTEST LABEL: \texttt{softmax\_online} in a 138-unit box.  At
%  \figunit = 0.13906mm that box is 54.4pt wide, and the label needs
%  roughly 59pt at \scriptsize.  If it overflows, either set \figMod to
%  \tiny or widen that column.
% =====================================================================

\providecommand{\figunit}{0.13906mm}
\providecommand{\figTitle}{\small}      % Host, DDR
\providecommand{\figMod}{\scriptsize}   % module names, section headings
\providecommand{\figNote}{\tiny}        % annotations and body text

\definecolor{figink}{HTML}{1A1A1A}
\definecolor{fighost}{HTML}{F2F2F2}
\definecolor{figddr}{HTML}{D8D8D8}
\definecolor{figlane}{HTML}{FAFAFA}
\definecolor{figcell}{HTML}{F0F0F0}
\definecolor{figrulec}{HTML}{9A9A9A}
\definecolor{figdim}{HTML}{555555}
\definecolor{figbody}{HTML}{333333}
\definecolor{figleg}{HTML}{666666}

\begin{tikzpicture}[
    x = {(\figunit,0mm)},
    y = {(0mm,\figunit)},
    line width = 0.5pt,
    draw       = figink,
    every node/.style = {inner sep=0pt, outer sep=0pt, text=figink},
    box/.style    = {draw=figink, fill=white},
    arw/.style    = {draw=figink, -{Stealth[length=3.2pt,width=2.6pt]}},
    rul/.style    = {draw=figrulec, line width=0.3pt},
    mod/.style    = {anchor=base, font=\figMod\ttfamily},
    ttl/.style    = {anchor=base, font=\figTitle},
    hed/.style    = {anchor=base west, font=\figMod},
    nte/.style    = {anchor=base, font=\figNote, text=figbody},
    dim/.style    = {anchor=base, font=\figNote, text=figdim},
    nteW/.style   = {anchor=base west, font=\figNote, text=figbody},
    dimW/.style   = {anchor=base west, font=\figNote, text=figdim},
]

% ---------------------------------------------------------------- regions
\draw[box, fill=fighost] (40,35)   rectangle (205,345);
\draw[box]               (235,35)  rectangle (1045,345);
\draw[box, fill=figddr]  (1075,35) rectangle (1250,345);

% ------------------------------------------------- write path, left to right
\draw[box] (250,222) rectangle (388,285);
\draw[box] (410,222) rectangle (548,285);
\draw[box] (570,222) rectangle (708,285);
\draw[box] (730,222) rectangle (868,285);
\draw[box] (890,222) rectangle (1028,285);

% --------------------- lane array; kv_store_ddr sits outside it, being shared
\draw[box, fill=figlane] (240,75) rectangle (872,163);
\draw[box] (250,88) rectangle (388,151);
\draw[box] (410,88) rectangle (548,151);
\draw[box] (570,88) rectangle (708,151);
\draw[box] (730,88) rectangle (868,151);

\draw[box] (890,77) rectangle (1028,161);
\draw[box, fill=figcell] (897,89) rectangle (927,117);
\draw[box, fill=figcell] (933,89) rectangle (963,117);
\draw[box, fill=figcell] (987,89) rectangle (1021,117);

% ---------------------------------------------------------------- arrows
\draw[arw] (208,254)  -- (246,254);
\draw[arw] (390,254)  -- (408,254);
\draw[arw] (550,254)  -- (568,254);
\draw[arw] (710,254)  -- (728,254);
\draw[arw] (870,254)  -- (888,254);
\draw[arw] (1032,254) -- (1071,254);

\draw[arw] (1071,119) -- (1032,119);
\draw[arw] (888,119)  -- (869,119);
\draw[arw] (728,119)  -- (710,119);
\draw[arw] (568,119)  -- (550,119);
\draw[arw] (408,119)  -- (390,119);
\draw[arw] (246,119)  -- (208,119);

\draw[arw] (479,222) -- (479,191) -- (799,191) -- (799,152);

% ------------------------------------------------- rules in host and DDR
\draw[rul] (56,301)   -- (189,301);
\draw[rul] (56,197)   -- (189,197);
\draw[rul] (56,123)   -- (189,123);
\draw[rul] (1090,285) -- (1235,285);
\draw[rul] (1090,235) -- (1235,235);
\draw[rul] (1090,143) -- (1235,143);

% ---------------------------------------------------------------- host
\node[ttl]  at (122,317) {Host};
\node[hed]  at (56,279)  {sent to the block};
\node[dimW] at (64,260)  {$q$, $k$ and $v$ for every};
\node[dimW] at (64,245)  {head, already projected};
\node[dimW] at (64,230)  {through $W_q$, $W_k$, $W_v$};
\node[dimW] at (64,215)  {and carrying RoPE};
\node[hed]  at (56,175)  {back from the block};
\node[dimW] at (64,156)  {the attention output};
\node[dimW] at (64,141)  {for all of the heads};
\node[hed]  at (56,101)  {finished on the host};
\node[dimW] at (64,82)   {$W_o'$ \ ($R^{-1}$ folded in)};
\node[dimW] at (64,67)   {MLP \ $\rightarrow$ \ logits};
\node[dimW] at (64,52)   {then the next layer};

% ---------------------------------------------------------------- DDR
\node[ttl] at (1162,319) {DDR};
\node[anchor=base, font=\figMod] at (1162,299) {KV cache};
\node[nte] at (1162,265) {one 52-byte row per};
\node[nte] at (1162,251) {token, per KV head};
\node[nte] at (1162,215) {the row is cut into $P$};
\node[nte] at (1162,201) {slices, and each one is};
\node[nte] at (1162,187) {kept as its own array,};
\node[nte] at (1162,173) {contiguous and running};
\node[nte] at (1162,159) {across all the tokens};
\node[nte] at (1162,123) {each AXI port reads};
\node[nte] at (1162,109) {one slice and nothing};
\node[nte] at (1162,95)  {else, which turns the};
\node[nte] at (1162,81)  {whole scan into long};
\node[nte] at (1162,67)  {sequential bursts};

% ------------------------------------------------------- programmable logic
\node[hed]  at (250,322) {Programmable logic};
\node[dimW] at (250,305) {one KV head at a time --- the $H$ heads of a token run one after the other};
\node[nteW] at (240,169) {$\times\,G$ \ --- one lane per query head of this KV head};

% ---------------------------------------------------------------- write path
\node[mod] at (319,249) {attn\_ingress};
\node[mod] at (479,258) {rot\_fwht};
\node[dim] at (479,239) {$k$, $v$, $q$};
\node[mod] at (639,249) {rot\_norm};
\node[mod] at (799,249) {rot\_encode};
\node[mod] at (959,262) {kv\_write};
\node[dim] at (959,245) {one row per token,};
\node[dim] at (959,233) {per KV head};

% ---------------------------------------------------------------- lanes
\node[mod] at (319,115) {finalize};
\node[mod] at (479,115) {accum};
\node[mod] at (639,115) {softmax\_online};
\node[mod] at (799,115) {score\_lane};

% ------------------------------------------------------- shared read engine
\node[mod] at (959,138) {kv\_store\_ddr};
\node[nte] at (912,98)  {0};
\node[nte] at (948,98)  {1};
\node[dim] at (975,97)  {$\dots$};
\node[nte] at (1004,98) {$P\!-\!1$};
\node[nte] at (959,64)  {the $P$ ports fetch one};
\node[nte] at (959,52)  {whole row per cycle, and};
\node[nte] at (959,40)  {all $G$ lanes score on it};

% ---------------------------------------------------------------- q branch
\node[nte] at (639,198) {$q$ is held inside each lane and never written to the cache};

% ---------------------------------------------------------------- legend
\node[anchor=base west, font=\figNote, text=figleg] at (40,15)
  {The ZCU104 provides four AXI-HP ports, each carrying 16 bytes a cycle. A plane never straddles a field boundary, so the key codes, the value codes and the norms must together fall into at most four planes};

\end{tikzpicture}

%        \caption{...}\label{fig:hw}
%      \end{figure*}
%
%  Preamble needs:
%      \usepackage{tikz}
%      \usetikzlibrary{arrows.meta}
%
%  NOT YET COMPILED -- written without a LaTeX toolchain available.
%  Check it renders before trusting it.
%
%  ---------------------------------------------------------------------
%  Geometry is in the same 1280 x 455 coordinate space as
%  fig_hardware.svg, so the two can be compared line for line.  The unit
%  vectors below set the printed size; text is set in real LaTeX sizes
%  and deliberately does NOT scale with the drawing.
%
%  To resize the drawing   : change \figunit
%  To resize the type      : change \figTitle / \figMod / \figNote
%
%  \figunit = 0.13906mm gives a width of 178mm (7in), a two-column
%  figure*.  For a single-column figure use about 0.0664mm, but at that
%  size the type will have to come down with it.
%
%  TIGHTEST LABEL: \texttt{softmax\_online} in a 138-unit box.  At
%  \figunit = 0.13906mm that box is 54.4pt wide, and the label needs
%  roughly 59pt at \scriptsize.  If it overflows, either set \figMod to
%  \tiny or widen that column.
% =====================================================================

\providecommand{\figunit}{0.13906mm}
\providecommand{\figTitle}{\small}      % Host, DDR
\providecommand{\figMod}{\scriptsize}   % module names, section headings
\providecommand{\figNote}{\tiny}        % annotations and body text

\definecolor{figink}{HTML}{1A1A1A}
\definecolor{fighost}{HTML}{F2F2F2}
\definecolor{figddr}{HTML}{D8D8D8}
\definecolor{figlane}{HTML}{FAFAFA}
\definecolor{figcell}{HTML}{F0F0F0}
\definecolor{figrulec}{HTML}{9A9A9A}
\definecolor{figdim}{HTML}{555555}
\definecolor{figbody}{HTML}{333333}
\definecolor{figleg}{HTML}{666666}

\begin{tikzpicture}[
    x = {(\figunit,0mm)},
    y = {(0mm,\figunit)},
    line width = 0.5pt,
    draw       = figink,
    every node/.style = {inner sep=0pt, outer sep=0pt, text=figink},
    box/.style    = {draw=figink, fill=white},
    arw/.style    = {draw=figink, -{Stealth[length=3.2pt,width=2.6pt]}},
    rul/.style    = {draw=figrulec, line width=0.3pt},
    mod/.style    = {anchor=base, font=\figMod\ttfamily},
    ttl/.style    = {anchor=base, font=\figTitle},
    hed/.style    = {anchor=base west, font=\figMod},
    nte/.style    = {anchor=base, font=\figNote, text=figbody},
    dim/.style    = {anchor=base, font=\figNote, text=figdim},
    nteW/.style   = {anchor=base west, font=\figNote, text=figbody},
    dimW/.style   = {anchor=base west, font=\figNote, text=figdim},
]

% ---------------------------------------------------------------- regions
\draw[box, fill=fighost] (40,35)   rectangle (205,345);
\draw[box]               (235,35)  rectangle (1045,345);
\draw[box, fill=figddr]  (1075,35) rectangle (1250,345);

% ------------------------------------------------- write path, left to right
\draw[box] (250,222) rectangle (388,285);
\draw[box] (410,222) rectangle (548,285);
\draw[box] (570,222) rectangle (708,285);
\draw[box] (730,222) rectangle (868,285);
\draw[box] (890,222) rectangle (1028,285);

% --------------------- lane array; kv_store_ddr sits outside it, being shared
\draw[box, fill=figlane] (240,75) rectangle (872,163);
\draw[box] (250,88) rectangle (388,151);
\draw[box] (410,88) rectangle (548,151);
\draw[box] (570,88) rectangle (708,151);
\draw[box] (730,88) rectangle (868,151);

\draw[box] (890,77) rectangle (1028,161);
\draw[box, fill=figcell] (897,89) rectangle (927,117);
\draw[box, fill=figcell] (933,89) rectangle (963,117);
\draw[box, fill=figcell] (987,89) rectangle (1021,117);

% ---------------------------------------------------------------- arrows
\draw[arw] (208,254)  -- (246,254);
\draw[arw] (390,254)  -- (408,254);
\draw[arw] (550,254)  -- (568,254);
\draw[arw] (710,254)  -- (728,254);
\draw[arw] (870,254)  -- (888,254);
\draw[arw] (1032,254) -- (1071,254);

\draw[arw] (1071,119) -- (1032,119);
\draw[arw] (888,119)  -- (869,119);
\draw[arw] (728,119)  -- (710,119);
\draw[arw] (568,119)  -- (550,119);
\draw[arw] (408,119)  -- (390,119);
\draw[arw] (246,119)  -- (208,119);

\draw[arw] (479,222) -- (479,191) -- (799,191) -- (799,152);

% ------------------------------------------------- rules in host and DDR
\draw[rul] (56,301)   -- (189,301);
\draw[rul] (56,197)   -- (189,197);
\draw[rul] (56,123)   -- (189,123);
\draw[rul] (1090,285) -- (1235,285);
\draw[rul] (1090,235) -- (1235,235);
\draw[rul] (1090,143) -- (1235,143);

% ---------------------------------------------------------------- host
\node[ttl]  at (122,317) {Host};
\node[hed]  at (56,279)  {sent to the block};
\node[dimW] at (64,260)  {$q$, $k$ and $v$ for every};
\node[dimW] at (64,245)  {head, already projected};
\node[dimW] at (64,230)  {through $W_q$, $W_k$, $W_v$};
\node[dimW] at (64,215)  {and carrying RoPE};
\node[hed]  at (56,175)  {back from the block};
\node[dimW] at (64,156)  {the attention output};
\node[dimW] at (64,141)  {for all of the heads};
\node[hed]  at (56,101)  {finished on the host};
\node[dimW] at (64,82)   {$W_o'$ \ ($R^{-1}$ folded in)};
\node[dimW] at (64,67)   {MLP \ $\rightarrow$ \ logits};
\node[dimW] at (64,52)   {then the next layer};

% ---------------------------------------------------------------- DDR
\node[ttl] at (1162,319) {DDR};
\node[anchor=base, font=\figMod] at (1162,299) {KV cache};
\node[nte] at (1162,265) {one 52-byte row per};
\node[nte] at (1162,251) {token, per KV head};
\node[nte] at (1162,215) {the row is cut into $P$};
\node[nte] at (1162,201) {slices, and each one is};
\node[nte] at (1162,187) {kept as its own array,};
\node[nte] at (1162,173) {contiguous and running};
\node[nte] at (1162,159) {across all the tokens};
\node[nte] at (1162,123) {each AXI port reads};
\node[nte] at (1162,109) {one slice and nothing};
\node[nte] at (1162,95)  {else, which turns the};
\node[nte] at (1162,81)  {whole scan into long};
\node[nte] at (1162,67)  {sequential bursts};

% ------------------------------------------------------- programmable logic
\node[hed]  at (250,322) {Programmable logic};
\node[dimW] at (250,305) {one KV head at a time --- the $H$ heads of a token run one after the other};
\node[nteW] at (240,169) {$\times\,G$ \ --- one lane per query head of this KV head};

% ---------------------------------------------------------------- write path
\node[mod] at (319,249) {attn\_ingress};
\node[mod] at (479,258) {rot\_fwht};
\node[dim] at (479,239) {$k$, $v$, $q$};
\node[mod] at (639,249) {rot\_norm};
\node[mod] at (799,249) {rot\_encode};
\node[mod] at (959,262) {kv\_write};
\node[dim] at (959,245) {one row per token,};
\node[dim] at (959,233) {per KV head};

% ---------------------------------------------------------------- lanes
\node[mod] at (319,115) {finalize};
\node[mod] at (479,115) {accum};
\node[mod] at (639,115) {softmax\_online};
\node[mod] at (799,115) {score\_lane};

% ------------------------------------------------------- shared read engine
\node[mod] at (959,138) {kv\_store\_ddr};
\node[nte] at (912,98)  {0};
\node[nte] at (948,98)  {1};
\node[dim] at (975,97)  {$\dots$};
\node[nte] at (1004,98) {$P\!-\!1$};
\node[nte] at (959,64)  {the $P$ ports fetch one};
\node[nte] at (959,52)  {whole row per cycle, and};
\node[nte] at (959,40)  {all $G$ lanes score on it};

% ---------------------------------------------------------------- q branch
\node[nte] at (639,198) {$q$ is held inside each lane and never written to the cache};

% ---------------------------------------------------------------- legend
\node[anchor=base west, font=\figNote, text=figleg] at (40,15)
  {The ZCU104 provides four AXI-HP ports, each carrying 16 bytes a cycle. A plane never straddles a field boundary, so the key codes, the value codes and the norms must together fall into at most four planes};

\end{tikzpicture}

  \caption{The accelerator. Attention is split once. Everything whose
  cost is constant in context length stays on the host, while the KV
  cache and the scan over it sit in programmable logic. Vermilion
  follows a token on its way into the cache and blue follows the scan on
  its way back out. The read engine sits outside the replicated lane
  array because one row fetch serves all $G$ query heads of a KV head.}
  \label{fig:hw}
\end{figure*}

\section{Architecture}
\label{sec:arch}

Figure~\ref{fig:hw} shows the organisation of the accelerator. The design
partitions attention at a single point. Every operation whose cost is
constant in context length remains on the host, while the KV cache and
the scan over it are placed in programmable logic. The block is invoked
once per layer and once per decode step. A single invocation processes
that layer's $H$ KV heads one after another, and within each head the
$G$ query heads that share it proceed in parallel. In the model
evaluated here $H{=}8$ and $G{=}4$, giving 32 attention heads of
dimension $d{=}64$.

\subsection{Partitioning}

The host computes the token embedding, the three projections $W_q$,
$W_k$ and $W_v$, and applies rotary position embedding. Position
therefore enters the datapath before the cache and never within it.

The target platform is a Xilinx ZCU104, whose Zynq UltraScale+ device
couples a processing system to programmable logic through AXI4
interfaces. Vectors cross to the accelerator over a 512-bit AXI4-Stream
port, each of the $d$ lanes carried as a 24-bit fixed-point value. Every
invocation receives $H$ groups of $2{+}G$ vectors, namely one key, one
value and $G$ queries for each KV head.

Keys and values are encoded and written to DDR before the scan begins,
because a token attends to itself as well as to every token before it.
The scan therefore covers $T$ rows, where $T$ counts the tokens seen so
far including the current one.

\subsection{Rotation and quantisation}

The rotation unit \texttt{rot\_fwht} applies one round of a randomised
Walsh--Hadamard rotation to keys, values and queries alike. All
subsequent arithmetic is performed in this rotated basis, and the
inverse rotation is folded into $W_o$ offline so that none is performed
at run time.

The norm unit \texttt{rot\_norm} computes the root-mean-square of the
rotated vector at the precision at which it will be stored. It is the
only producer of this norm, which guarantees that the quantiser's
decision thresholds and the eventual decoder operate on the same value.

The encoder \texttt{rot\_encode} maps each rotated vector to codewords
without division. Comparing a normalised channel against a threshold is
equivalent to comparing the raw channel against that threshold scaled by
the norm, so the thresholds are scaled once per vector and each of the
$d$ channels reduces to a comparison. Separate instances serve keys and
values, which differ in codeword width and in threshold table.

\subsection{The cache row}

For one token and one KV head the cache holds $d$ key codewords, $d$
value codewords, and two scalar norms, one for the key vector and one
for the value vector. Writing $b_k$ and $b_v$ for the codeword widths
and $w_n$ for the norm width, a row occupies

\begin{equation}
  R \;=\;
  \underbrace{d\,b_k}_{\text{key codewords}} \;+\;
  \underbrace{d\,b_v}_{\text{value codewords}} \;+\;
  \underbrace{2\,w_n}_{\text{two norms}}
  \quad \text{bits.}
  \label{eq:row}
\end{equation}

In the configuration built, $d{=}64$, $b_k{=}4$, $b_v{=}2$ and
$w_n{=}16$, so a row is 416 bits, or 52 bytes. The write unit
\texttt{kv\_write} assembles this row and issues it to DDR.

\subsection{Cache layout}

The cache resides entirely in DDR and is stored transposed. A row is
partitioned into $P$ slices, and each slice is held as a contiguous
array spanning all tokens.

This layout follows from a bandwidth constraint. The scan consumes one
row per cycle, while a single 128-bit AXI high-performance port supplies
16 bytes per cycle. $P$ ports must therefore operate concurrently, and
each requires its own sequential access stream. A row-major layout
contains only one such stream. Partitioning the token range across ports
does not help either, because the causal scan advances through tokens in
order and would leave all but one port idle. Under the transposed layout
each port traverses a single slice, and the $P$ streams jointly deliver
one complete row every cycle.

A slice never spans a field boundary, so the port requirement is not a
function of row size alone. The key codewords, the value codewords and
the norms must together occupy at most four slices, which is the number
of high-performance ports the device provides. At 250\,MHz the scan
sustains 1.14 cycles per row against a one-cycle ideal.

\subsection{Scan datapath}

Attention is computed directly on codewords rather than on reconstructed
vectors. A cached key is the product of its norm and a codebook entry
selected by its codeword, so the inner product with a query is a sum
over channels of query element times codebook entry, scaled once by that
norm. The per-channel product can therefore take only $2^{b_k}$ distinct
values, one for each codebook entry.

The table-building unit \texttt{qtab\_build} evaluates those $2^{b_k}$
products for every channel once per query vector. Each cached token
afterwards costs $d$ table lookups, an adder tree, and a single
multiplication by the row norm, which is two DSP slices per lane rather
than $d$.

The read engine \texttt{kv\_store\_ddr} is instantiated once per block
and shared by all $G$ lanes, so one row fetch serves $G$ query heads.
This makes the reuse already present in grouped-query attention explicit
in the hardware.

Each lane then scores, normalises and accumulates. The scoring unit
\texttt{score\_lane} produces the score for the row. The softmax unit
\texttt{softmax\_online} maintains a running maximum and rescales the
accumulator whenever that maximum grows, which allows the scan to finish
in a single pass. A two-pass formulation would require a second
traversal of an off-die cache and is not affordable. A small lookup
table \texttt{exp\_lut} evaluates the exponential by splitting it into a
shift and a table index, so the table need only cover the unit interval
however large its argument becomes. The accumulator \texttt{accum}
weights the decoded value and adds it to a running sum, exploiting the
same property on the value side, since with $2^{b_v}$ value centroids
one multiplication produces a shared factor and every channel reduces to
a multiplexer and an adder. The normalisation unit \texttt{finalize}
inverts the softmax denominator once per head and multiplies across
channels, the denominator not being final until the scan ends.

\subsection{Egress}

Each group returns its $G$ outputs over the same 512-bit stream, after
which the block advances to the next KV head. The host then applies
$W_o$, followed by the MLP and the residual connection. A decode step
completes in 22 to 26 microseconds on the device.

Prompt tokens traverse this path identically, so every row the scan
reads was written by the accelerator itself.
